\documentclass[acmsmall,screen,review,anonymous]{acmart}

\setcopyright{none}
\acmConference[FSE 2027]{FSE 2027 Ideas, Visions and Reflections}{February 2027}{Shenzhen, China}
\acmBooktitle{FSE 2027: Ideas, Visions and Reflections}

\begin{document}

\title[Registries That Never Look Again]{Registries That Never Look Again: Stale Install Instructions in an Agent Component Registry}

\begin{abstract}
Agent component registries tell a client \emph{what to install}. We measured whether that
instruction still works. On one production registry's snapshot of 2026-11-16 (65{,}340
latest records, 22{,}526 distinct npm/PyPI references) 274 references named a package or
version that did not resolve; after correcting one normalisation defect the figure is at
most 267 (1.19\%). Running the actual install command for 40 of them failed 39 times
(97.5\%, Wilson 95\% CI 87.1--99.6\%). The entries carrying those references are not fresh
mistakes: their median age is 218 days (p90 283, max 435) and 83\% are older than 90 days.
The packages are not abandoned either --- of the 140 whose download statistics were
obtainable, the median is 49.5 downloads per month and 47 exceed 100. A registry that
validates only at publication therefore keeps advertising unusable installs, for months,
for packages people use. Our counts of broken references are \emph{lower bounds}: those rejected at
publish time were never published and are invisible to us.
\end{abstract}

\keywords{software registries, package installation, staleness, measurement}

\maketitle

\section{The Claim}
A registry entry is an instruction. This paper asks three questions about that instruction:
does it work, how long has it not worked, and is anyone following it? Table~\ref{tab:rq}
gives the answers. The counts in the first two rows are lower bounds, because we can only
observe references that were published at all.

\begin{table}[t]
\caption{Three questions, three measurements, one snapshot (2026-11-16).}
\label{tab:rq}
\small
\begin{tabular}{@{}lll@{}}
\toprule
Question & Measurement & Result\\
\midrule
Does the install work? & run \texttt{npm install} / \texttt{pip download} for 40 & 39/40 fail (97.5\%)\\
How long dead? & entry age from the registry's own timestamp & median 218 d\\
Does anyone use it? & monthly downloads per package & median 49.5/mo; 47 $\geq$ 100\\
\bottomrule
\end{tabular}
\end{table}

\section{Corpus}
The corpus is the full pagination of the registry's \texttt{?version=latest} endpoint on
2026-11-16: 654 pages, 65{,}340 latest records, de-duplicated to 22{,}526 distinct
\texttt{(registry, identifier, version)} references across 22{,}384 servers. Each reference
was resolved against the publisher's own registry: the npm \texttt{versions} map and the
PyPI \texttt{releases} map, requiring at least one file. 274 references did not resolve.

\paragraph{Two corrections we made mid-study.} Both came from re-running the real thing.
\emph{(i) PEP~440}: PyPI keys \texttt{releases} by normal form~\cite{pep440} and we compared
the raw string, so seven of the 32 PyPI version-absent rows were not defects.
\emph{(ii) npm tombstones}: an unpublished unscoped package still answers \texttt{200} with
an empty \texttt{versions} map and a \texttt{time.unpublished} object, which we read as
``package exists'', filing 68 tombstones as \emph{version} absent. Corrected:
\textbf{267} unresolvable --- \textbf{219} gone, \textbf{48} version-absent (1.19\%). The
total barely moved; the split did, and with it the story: wholesale removal, not drift.

\section{Does the Install Work?}
We sampled 20 rows from each failure class at equal spacing over the sorted list --- not
hand-picked --- and executed:
\begin{quote}\small
\texttt{npm install <name>@<ver> --dry-run --ignore-scripts}\\
\texttt{python -m pip download <name>==<ver> --no-deps}
\end{quote}
Table~\ref{tab:repro} reports the outcome. Both failure modes are covered, and they fail
with different error classes: npm \texttt{404} and \texttt{ETARGET notarget}, pip
\texttt{No matching distribution}. The single non-failing case is the PEP~440 case above,
so it is evidence for the correction rather than against the finding.

\begin{table}[t]
\caption{Real install attempts, equal-spaced sample of the published list.}
\label{tab:repro}
\small
\begin{tabular}{@{}lrrrl@{}}
\toprule
class & $n$ & reproduced & rate & 95\% Wilson CI\\
\midrule
package-gone & 20 & 20 & 100\% & 83.9--100\%\\
version-absent & 20 & 19 & 95\% & 76.4--99.1\%\\
\midrule
all & 40 & 39 & 97.5\% & 87.1--99.6\%\\
\bottomrule
\end{tabular}
\end{table}

\paragraph{They are not transient.} Re-resolving all 274 references three days later gives
267 still broken and seven ``recovered'' --- and all seven are exactly the PEP~440 cases.
There were \emph{zero} upstream repairs in three days.

\section{How Long Has It Been Dead?}
Joining the 274 references back to the snapshot yields 311 entries (a reference can appear
under several servers). Reading each entry's own
\texttt{\_meta[...].publishedAt}: minimum 2.0 days, p25 128, \textbf{median 218} (7.2
months), p75 266, p90 283, maximum 435 days; 96.8\% are older than 30 days and 83.3\%
older than 90. Split by class, package-gone has a median of 245 days and version-absent
192.

\paragraph{How to read it.} \texttt{publishedAt} is when the registry published the entry,
which cannot precede the upstream disappearance. Entry age is therefore an \emph{upper
bound} on the break duration, not the duration itself: the claim is that the registry kept
a dead reference for at least a median of seven months, not that the install has been
broken for seven months.

\section{Does Anyone Use It?}
Monthly downloads for the 271 distinct packages, from the npm downloads API and PyPI
download statistics, throttled after rate limits: 140 packages returned statistics with a
median of 49.5 downloads per month (p90 666, max 20{,}925), and 47 packages exceed 100 per
month. A further 48 packages return 404 on the download API, which is itself consistent
with ``gone''.

\paragraph{How to read it.} Downloads are per \emph{package}, not per \emph{version}. The
20{,}925 downloads of the busiest package are overwhelmingly other versions. The claim is
that the package is in active use and the exact version the registry advertises cannot be
installed --- not that twenty thousand users are affected.

\section{Why This Happens, and How General It Is}
\paragraph{The general form.} Any registry that validates a reference \emph{only at
publication} accumulates references that later become unexecutable, at whatever rate
upstream withdrawals occur. Without re-validation the residence time of a stale reference
is unbounded, so by Little's law the stock grows monotonically: registry accuracy is a
function of the re-check interval, not of publish-time rigour.

\paragraph{This registry is a favourable instance.} It does validate at publish and edit
time and timestamps every entry, so 1.19\% is what staleness looks like \emph{when the
publisher is careful}. A registry that validates less would show more, and the references
rejected outright never appear in our count.

\paragraph{The cheap fix.} Re-run the existing validation on a schedule and record the
outcome in metadata (\texttt{lastVerifiedAt}, \texttt{verificationStatus} $\in$ \{ok,
package-missing, version-missing\}), leaving the entry's history intact. Deletion is a
separate decision: revocation across distributions is not the registry's call alone.

\section{Threats to Validity}
\textbf{Sample size.} RQ1 uses 40 of 274 references; the interval is reported and the point
estimate should not be read as an exact population rate. \textbf{Upper bound on age} and
\textbf{package-level downloads} are as discussed above. \textbf{Snapshot drift}: 274 on
2026-11-16 became 275 on 2026-11-17 under the same definition. \textbf{Short re-check
window}: three days shows these are not transient but cannot show annual inflow; a second
full audit on 2026-12-15 will report the new count, the recovery count and the weekly inflow
of new breakage. \textbf{Residual normalisation risk} in npm version strings was not
exhaustively ruled out.

\section{Related Work}
Two adjacent lines study agent supply chains. One maps the risks and the open problems of the
LLM supply chain as a whole~\cite{llmsupplychain}; the other certifies a structured action
against several evidence sources under a corruption budget~\cite{multisource}. Both assume
the inputs they receive are the current ones. This paper asks the prior question --- whether
the registry's own instruction still resolves --- and finds a failure with no adversary and
no intent, only a stale pointer.

The distribution formats themselves are specified: PEP~440 defines how Python versions are
normalised (the defect this paper had to correct)~\cite{pep440}, and PEP~740 defines how
publishers attach attestations to release files~\cite{pep740}. Neither specifies a
re-validation duty for an index or a registry.

\section{Conclusion}
A registry that validates once and never again accumulates instructions that no longer
work. In one production snapshot that accumulation is already measurable: 267 unresolved
references at most, 39 of 40 sampled installs failing when actually executed, a median
entry age of seven months, and 47 affected packages still receiving over a hundred
downloads a month. The fix is cheap and does not require deleting history.


{\footnotesize
\begin{thebibliography}{9}
\bibitem{pep440} Python Software Foundation.
\newblock {PEP} 440: Version Identification and Dependency Specification.
\newblock \url{https://peps.python.org/pep-0440/}

\bibitem{pep740} Python Software Foundation.
\newblock {PEP} 740: Index support for digital attestations.
\newblock \url{https://peps.python.org/pep-0740/}

\bibitem{llmsupplychain} Q.~Hu, X.~Xie, S.~Chen, and L.~Ma.
\newblock Large Language Model Supply Chain: Open Problems From the Security Perspective.
\newblock arXiv:2411.01604, 2024.

\bibitem{multisource} A.~Pandey, A.~Jain, L.~Chen, C.~Maple, and C.~Panchev.
\newblock Certified Multi-Source Integrity for Structured Agent Actions.
\newblock arXiv:2609.34245, 2026.

\bibitem{registry} Model Context Protocol.
\newblock Registry specification and server metadata (\texttt{server.json}).
\newblock \url{https://github.com/modelcontextprotocol/registry}
\end{thebibliography}
}

\end{document}
